% --- Archon physics formalization source begin ---
% archon:physics
% archon:covers PhyXMiniProblems/problem_phyx_mini_0064.lean
% archon:source-report reports/phyx_mini/problem_phyx_mini_0064.source.json

\chapter{Physics problem phyx\_mini\_0064}
\label{ch:PhyXMiniProblems_problem_phyx_mini_0064}

\paragraph{Problem source.}
\#\# Question

Find the focal length of the glass lens in Figure

\#\# Answer choices

- A: A: 24cm

- B: B: 56cm

- C: C: 30cm

- D: D: 35cm

\#\# Figure caption (auxiliary; use the image as primary evidence)

The image shows a diagram of a lens with two labeled distances. A convex lens is illustrated in the center, represented by a symmetrical shape shaded in blue. A horizontal line passes through the lens, indicating the principal axis.

- On the left, there's a black dot along the axis with a line extending to the center of the lens, labeled "40 cm."
- On the right of the lens, another black dot is present with a similar line extending towards the center, labeled "24 cm."

The two lines with measurements suggest object and image distances relative to the lens.

\#\# Dataset metadata

Geometrical Optics; Physical Model Grounding Reasoning, Multi-Formula Reasoning

\paragraph{Recorded answer/context.}
C: C: 30cm

\paragraph{Figure/image path.}
/root/proposal\_for\_physic/science-mango/phyx\_mini\_run/phyx\_data/test\_image/64.png

\paragraph{Formalization target.}
create a compiling Lean file with sorry bodies at `PhyXMiniProblems/problem\_phyx\_mini\_0064.lean`. The Lean declarations must preserve the physical quantities, dimensions or dimensional roles, figure labels, governing-law hypotheses, and final relation expressed by this problem.
Use Mathlib/Physlib names found through LeanExplore where available. If a physics API is missing, introduce faithful local abstractions rather than scalar placeholder aliases.

\begin{theorem}[Physics formalization target]
\label{thm:physics:phyx_mini_0064:target}
\lean{PhyXMiniProblems.ProblemPhyXMini0064.problem_phyx_mini_0064}
\uses{def:physics:phyx-mini-0064:phyxminiproblems-problemphyxmini0064-lengthincentimeters, def:physics:phyx-mini-0064:phyxminiproblems-problemphyxmini0064-glasslensmakersetup, def:physics:phyx-mini-0064:phyxminiproblems-problemphyxmini0064-matchesfiguregeometryandreadouts, def:physics:phyx-mini-0064:phyxminiproblems-problemphyxmini0064-usesstandardglassinairindices, def:physics:phyx-mini-0064:phyxminiproblems-problemphyxmini0064-hasphysicallensparameters, def:physics:phyx-mini-0064:phyxminiproblems-problemphyxmini0064-satisfiesthinlensmakerequation, def:physics:phyx-mini-0064:phyxminiproblems-problemphyxmini0064-answerfocallengthincentimeters}
The assigned autoformalize agent should translate this physics problem into Lean declarations in the covered file, with theorem and lemma proof bodies written as `by sorry`.
\end{theorem}
\begin{proof}
This is an autoformalization task, not a proof task. Produce statements that can later be proved by the physics prover without weakening the physical model.
\end{proof}

% --- Archon declaration topology begin ---
\section{Declaration topology}
\begin{definition}[Length Magnitude]
  \label{def:physics:phyx-mini-0064:phyxminiproblems-problemphyxmini0064-lengthmagnitude}
  \lean{PhyXMiniProblems.ProblemPhyXMini0064.LengthMagnitude}
  A nonnegative physical quantity carrying the dimension of length
\end{definition}
\begin{proof}
  This is the declared model definition of Length Magnitude.
\end{proof}
\begin{definition}[length In Centimeters]
  \label{def:physics:phyx-mini-0064:phyxminiproblems-problemphyxmini0064-lengthincentimeters}
  \lean{PhyXMiniProblems.ProblemPhyXMini0064.lengthInCentimeters}
  \uses{def:physics:phyx-mini-0064:phyxminiproblems-problemphyxmini0064-lengthmagnitude}
  The numerical readout, in centimeters, of a physical length magnitude
\end{definition}
\begin{proof}
  This is the declared model definition of length In Centimeters.
\end{proof}
\begin{definition}[Axial Side]
  \label{def:physics:phyx-mini-0064:phyxminiproblems-problemphyxmini0064-axialside}
  \lean{PhyXMiniProblems.ProblemPhyXMini0064.AxialSide}
  Axial sides relative to light propagation from left to right
\end{definition}
\begin{proof}
  This is the declared model definition of Axial Side.
\end{proof}
\begin{definition}[Lens Surface Label]
  \label{def:physics:phyx-mini-0064:phyxminiproblems-problemphyxmini0064-lenssurfacelabel}
  \lean{PhyXMiniProblems.ProblemPhyXMini0064.LensSurfaceLabel}
  The two spherical faces of the lens, ordered along the propagation axis
\end{definition}
\begin{proof}
  This is the declared model definition of Lens Surface Label.
\end{proof}
\begin{definition}[Center Of Curvature Dot]
  \label{def:physics:phyx-mini-0064:phyxminiproblems-problemphyxmini0064-centerofcurvaturedot}
  \lean{PhyXMiniProblems.ProblemPhyXMini0064.CenterOfCurvatureDot}
  The two black centers of curvature drawn on the principal axis
\end{definition}
\begin{proof}
  This is the declared model definition of Center Of Curvature Dot.
\end{proof}
\begin{definition}[Optical Medium]
  \label{def:physics:phyx-mini-0064:phyxminiproblems-problemphyxmini0064-opticalmedium}
  \lean{PhyXMiniProblems.ProblemPhyXMini0064.OpticalMedium}
  The optical media on the two sides of each refracting interface
\end{definition}
\begin{proof}
  This is the declared model definition of Optical Medium.
\end{proof}
\begin{definition}[Propagation Direction]
  \label{def:physics:phyx-mini-0064:phyxminiproblems-problemphyxmini0064-propagationdirection}
  \lean{PhyXMiniProblems.ProblemPhyXMini0064.PropagationDirection}
  Direction used to assign incident/outgoing faces and signed radii
\end{definition}
\begin{proof}
  This is the declared model definition of Propagation Direction.
\end{proof}
\begin{definition}[Spherical Lens Surface]
  \label{def:physics:phyx-mini-0064:phyxminiproblems-problemphyxmini0064-sphericallenssurface}
  \lean{PhyXMiniProblems.ProblemPhyXMini0064.SphericalLensSurface}
  \uses{def:physics:phyx-mini-0064:phyxminiproblems-problemphyxmini0064-lengthmagnitude, def:physics:phyx-mini-0064:phyxminiproblems-problemphyxmini0064-axialside}
  This structure records the Spherical Lens Surface component of the problem-specific physical model
\end{definition}
\begin{proof}
  This is the declared model definition of Spherical Lens Surface.
\end{proof}
\begin{definition}[Thin Biconvex Glass Lens]
  \label{def:physics:phyx-mini-0064:phyxminiproblems-problemphyxmini0064-thinbiconvexglasslens}
  \lean{PhyXMiniProblems.ProblemPhyXMini0064.ThinBiconvexGlassLens}
  \uses{def:physics:phyx-mini-0064:phyxminiproblems-problemphyxmini0064-lengthmagnitude, def:physics:phyx-mini-0064:phyxminiproblems-problemphyxmini0064-lenssurfacelabel, def:physics:phyx-mini-0064:phyxminiproblems-problemphyxmini0064-sphericallenssurface}
  A physical thin biconvex lens with two labeled spherical faces
\end{definition}
\begin{proof}
  This is the declared model definition of Thin Biconvex Glass Lens.
\end{proof}
\begin{definition}[Curvature Radius Figure]
  \label{def:physics:phyx-mini-0064:phyxminiproblems-problemphyxmini0064-curvatureradiusfigure}
  \lean{PhyXMiniProblems.ProblemPhyXMini0064.CurvatureRadiusFigure}
  \uses{def:physics:phyx-mini-0064:phyxminiproblems-problemphyxmini0064-lengthmagnitude, def:physics:phyx-mini-0064:phyxminiproblems-problemphyxmini0064-axialside, def:physics:phyx-mini-0064:phyxminiproblems-problemphyxmini0064-lenssurfacelabel, def:physics:phyx-mini-0064:phyxminiproblems-problemphyxmini0064-centerofcurvaturedot}
  This structure records the Curvature Radius Figure component of the problem-specific physical model
\end{definition}
\begin{proof}
  This is the declared model definition of Curvature Radius Figure.
\end{proof}
\begin{definition}[Glass Lensmaker Setup]
  \label{def:physics:phyx-mini-0064:phyxminiproblems-problemphyxmini0064-glasslensmakersetup}
  \lean{PhyXMiniProblems.ProblemPhyXMini0064.GlassLensmakerSetup}
  \uses{def:physics:phyx-mini-0064:phyxminiproblems-problemphyxmini0064-opticalmedium, def:physics:phyx-mini-0064:phyxminiproblems-problemphyxmini0064-propagationdirection, def:physics:phyx-mini-0064:phyxminiproblems-problemphyxmini0064-thinbiconvexglasslens, def:physics:phyx-mini-0064:phyxminiproblems-problemphyxmini0064-curvatureradiusfigure}
  The physical lens, material data, propagation convention, and figure
\end{definition}
\begin{proof}
  This is the declared model definition of Glass Lensmaker Setup.
\end{proof}
\begin{definition}[signed Surface Radius In Centimeters]
  \label{def:physics:phyx-mini-0064:phyxminiproblems-problemphyxmini0064-signedsurfaceradiusincentimeters}
  \lean{PhyXMiniProblems.ProblemPhyXMini0064.signedSurfaceRadiusInCentimeters}
  \uses{def:physics:phyx-mini-0064:phyxminiproblems-problemphyxmini0064-lengthincentimeters, def:physics:phyx-mini-0064:phyxminiproblems-problemphyxmini0064-sphericallenssurface}
  This def records the signed Surface Radius In Centimeters component of the problem-specific physical model
\end{definition}
\begin{proof}
  This is the declared model definition of signed Surface Radius In Centimeters.
\end{proof}
\begin{definition}[Matches Figure Geometry And Readouts]
  \label{def:physics:phyx-mini-0064:phyxminiproblems-problemphyxmini0064-matchesfiguregeometryandreadouts}
  \lean{PhyXMiniProblems.ProblemPhyXMini0064.MatchesFigureGeometryAndReadouts}
  \uses{def:physics:phyx-mini-0064:phyxminiproblems-problemphyxmini0064-lengthincentimeters, def:physics:phyx-mini-0064:phyxminiproblems-problemphyxmini0064-glasslensmakersetup}
  This def records the Matches Figure Geometry And Readouts component of the problem-specific physical model
\end{definition}
\begin{proof}
  This is the declared model definition of Matches Figure Geometry And Readouts.
\end{proof}
\begin{definition}[Uses Standard Glass In Air Indices]
  \label{def:physics:phyx-mini-0064:phyxminiproblems-problemphyxmini0064-usesstandardglassinairindices}
  \lean{PhyXMiniProblems.ProblemPhyXMini0064.UsesStandardGlassInAirIndices}
  \uses{def:physics:phyx-mini-0064:phyxminiproblems-problemphyxmini0064-glasslensmakersetup}
  This def records the Uses Standard Glass In Air Indices component of the problem-specific physical model
\end{definition}
\begin{proof}
  This is the declared model definition of Uses Standard Glass In Air Indices.
\end{proof}
\begin{definition}[Has Physical Lens Parameters]
  \label{def:physics:phyx-mini-0064:phyxminiproblems-problemphyxmini0064-hasphysicallensparameters}
  \lean{PhyXMiniProblems.ProblemPhyXMini0064.HasPhysicalLensParameters}
  \uses{def:physics:phyx-mini-0064:phyxminiproblems-problemphyxmini0064-lengthincentimeters, def:physics:phyx-mini-0064:phyxminiproblems-problemphyxmini0064-glasslensmakersetup}
  Positivity and optical-branch conditions for the modeled lens
\end{definition}
\begin{proof}
  This is the declared model definition of Has Physical Lens Parameters.
\end{proof}
\begin{definition}[Satisfies Thin Lensmaker Equation]
  \label{def:physics:phyx-mini-0064:phyxminiproblems-problemphyxmini0064-satisfiesthinlensmakerequation}
  \lean{PhyXMiniProblems.ProblemPhyXMini0064.SatisfiesThinLensmakerEquation}
  \uses{def:physics:phyx-mini-0064:phyxminiproblems-problemphyxmini0064-lengthincentimeters, def:physics:phyx-mini-0064:phyxminiproblems-problemphyxmini0064-glasslensmakersetup, def:physics:phyx-mini-0064:phyxminiproblems-problemphyxmini0064-signedsurfaceradiusincentimeters}
  This def records the Satisfies Thin Lensmaker Equation component of the problem-specific physical model
\end{definition}
\begin{proof}
  This is the declared model definition of Satisfies Thin Lensmaker Equation.
\end{proof}
\begin{definition}[Answer Choice]
  \label{def:physics:phyx-mini-0064:phyxminiproblems-problemphyxmini0064-answerchoice}
  \lean{PhyXMiniProblems.ProblemPhyXMini0064.AnswerChoice}
  Labels of the four displayed multiple-choice answers
\end{definition}
\begin{proof}
  This is the declared model definition of Answer Choice.
\end{proof}
\begin{definition}[answer Focal Length In Centimeters]
  \label{def:physics:phyx-mini-0064:phyxminiproblems-problemphyxmini0064-answerfocallengthincentimeters}
  \lean{PhyXMiniProblems.ProblemPhyXMini0064.answerFocalLengthInCentimeters}
  \uses{def:physics:phyx-mini-0064:phyxminiproblems-problemphyxmini0064-answerchoice}
  Focal-length readout printed beside each answer choice, in centimeters
\end{definition}
\begin{proof}
  This is the declared model definition of answer Focal Length In Centimeters.
\end{proof}
\begin{lemma}[signed surface radius readouts]
  \label{lem:physics:phyx-mini-0064:phyxminiproblems-problemphyxmini0064-signed-surface-radius-readouts}
  \lean{PhyXMiniProblems.ProblemPhyXMini0064.signed_surface_radius_readouts}
  \uses{def:physics:phyx-mini-0064:phyxminiproblems-problemphyxmini0064-glasslensmakersetup, def:physics:phyx-mini-0064:phyxminiproblems-problemphyxmini0064-signedsurfaceradiusincentimeters, def:physics:phyx-mini-0064:phyxminiproblems-problemphyxmini0064-matchesfiguregeometryandreadouts}
  This lemma records the signed surface radius readouts component of the problem-specific physical model
\end{lemma}
\begin{proof}
  The conclusion follows from the governing relations and figure readouts named in the statement of signed surface radius readouts.
\end{proof}
\begin{lemma}[reciprocal focal length eq one over thirty]
  \label{lem:physics:phyx-mini-0064:phyxminiproblems-problemphyxmini0064-reciprocal-focal-length-eq-one-over-thirty}
  \lean{PhyXMiniProblems.ProblemPhyXMini0064.reciprocal_focal_length_eq_one_over_thirty}
  \uses{def:physics:phyx-mini-0064:phyxminiproblems-problemphyxmini0064-lengthincentimeters, def:physics:phyx-mini-0064:phyxminiproblems-problemphyxmini0064-glasslensmakersetup, def:physics:phyx-mini-0064:phyxminiproblems-problemphyxmini0064-matchesfiguregeometryandreadouts, def:physics:phyx-mini-0064:phyxminiproblems-problemphyxmini0064-usesstandardglassinairindices, def:physics:phyx-mini-0064:phyxminiproblems-problemphyxmini0064-satisfiesthinlensmakerequation}
  This lemma records the reciprocal focal length eq one over thirty component of the problem-specific physical model
\end{lemma}
\begin{proof}
  The conclusion follows from the governing relations and figure readouts named in the statement of reciprocal focal length eq one over thirty.
\end{proof}
% --- Archon declaration topology end ---
% --- Archon physics formalization source end ---
